\hypertarget{roadsos-structured-database-documentation}{%
\section{RoadSOS --- Structured Database
Documentation}\label{roadsos-structured-database-documentation}}

\textbf{Submission artifact for Section 7.1 of the IIT Madras Road
Safety Hackathon 2026 rulebook:} \textgreater{} \emph{``Structured
database used for the models are to be submitted.''}

This document describes every persistent data structure used by RoadSOS,
its schema, its source, and how it is consumed by the runtime.

\begin{center}\rule{0.5\linewidth}{0.5pt}\end{center}

\hypertarget{emergency-numbers-database}{%
\subsection{1. Emergency Numbers
Database}\label{emergency-numbers-database}}

\textbf{Authoritative source:}
\texttt{backend/data/emergency\_seed.json} \textbf{Frontend bundle
(generated):} \texttt{frontend/src/utils/emergencyNumbers.js}
\textbf{Endpoint exposing it:} \texttt{GET\ /offline-pack}

\hypertarget{coverage}{%
\subsubsection{1.1 Coverage}\label{coverage}}

\begin{longtable}[]{@{}ll@{}}
\toprule\noalign{}
Metric & Value \\
\midrule\noalign{}
\endhead
\bottomrule\noalign{}
\endlastfoot
Total countries \& territories & \textbf{200} \\
UN-recognised states covered & 100\% \\
Asia & 49 \\
Europe & 45 \\
Africa & 55 \\
Americas & 35 \\
Oceania & 16 \\
File size (raw JSON) & \textasciitilde28 KB \\
File size (gzipped) & \textasciitilde9 KB \\
\end{longtable}

\hypertarget{schema}{%
\subsubsection{1.2 Schema}\label{schema}}

\begin{verbatim}
{
  "country":       "string  — Human-readable English country name",
  "country_code":  "string  — ISO 3166-1 alpha-2 code, uppercase",
  "calling_code":  "string  — International dialling prefix incl. '+'",
  "police":        "string  — National police emergency number",
  "ambulance":     "string  — National ambulance / medical emergency number",
  "fire":          "string  — National fire department emergency number",
  "general":       "string  — Universal emergency number (often 112 / 911)"
}
\end{verbatim}

\hypertarget{example-record}{%
\subsubsection{1.3 Example record}\label{example-record}}

\begin{verbatim}
{
  "country": "India",
  "country_code": "IN",
  "calling_code": "+91",
  "police": "100",
  "ambulance": "108",
  "fire": "101",
  "general": "112"
}
\end{verbatim}

\hypertarget{data-sourcing-verification}{%
\subsubsection{1.4 Data sourcing \&
verification}\label{data-sourcing-verification}}

Primary source: \textbf{Wikipedia ``List of emergency telephone
numbers''} (canonical, government-maintained). Each entry was
cross-checked against: 1. The country's official government emergency
services portal where available 2. ITU-T E.164 numbering plan
documentation 3. EU-wide 112 standard for European jurisdictions

Where a country uses different numbers for different services, the
most-publicised national number is recorded. For countries where 112 is
universally accepted as a fallback (e.g.~all GSM networks),
\texttt{general} is set to \texttt{112}.

\hypertarget{update-policy}{%
\subsubsection{1.5 Update policy}\label{update-policy}}

Static. Emergency numbers change roughly once per decade. Changes are
made by editing \texttt{backend/data/emergency\_seed.json} and
regenerating the frontend bundle:

\begin{verbatim}
# Regenerate frontend/src/utils/emergencyNumbers.js from backend JSON
python scripts/sync_emergency_numbers.py
\end{verbatim}

\begin{center}\rule{0.5\linewidth}{0.5pt}\end{center}

\hypertarget{osm-category-classification-map}{%
\subsection{2. OSM Category Classification
Map}\label{osm-category-classification-map}}

\textbf{Source:} \texttt{backend/services/overpass\_service.py} →
\texttt{CATEGORY\_MAP} \textbf{Purpose:} Maps OpenStreetMap tag (key,
value) pairs to the six visible RoadSOS contact categories.

\hypertarget{schema-1}{%
\subsubsection{2.1 Schema}\label{schema-1}}

\begin{verbatim}
CATEGORY_MAP: list[tuple[tuple[str, str], str]]
# [ ((osm_key, osm_value), roadsos_category), ... ]
\end{verbatim}

\hypertarget{full-mapping}{%
\subsubsection{2.2 Full mapping}\label{full-mapping}}

\begin{longtable}[]{@{}
  >{\raggedright\arraybackslash}p{(\columnwidth - 6\tabcolsep) * \real{0.2500}}
  >{\raggedright\arraybackslash}p{(\columnwidth - 6\tabcolsep) * \real{0.2500}}
  >{\raggedright\arraybackslash}p{(\columnwidth - 6\tabcolsep) * \real{0.2500}}
  >{\raggedright\arraybackslash}p{(\columnwidth - 6\tabcolsep) * \real{0.2500}}@{}}
\toprule\noalign{}
\begin{minipage}[b]{\linewidth}\raggedright
OSM Key
\end{minipage} & \begin{minipage}[b]{\linewidth}\raggedright
OSM Value
\end{minipage} & \begin{minipage}[b]{\linewidth}\raggedright
RoadSOS Category
\end{minipage} & \begin{minipage}[b]{\linewidth}\raggedright
Notes
\end{minipage} \\
\midrule\noalign{}
\endhead
\bottomrule\noalign{}
\endlastfoot
\texttt{amenity} & \texttt{hospital} & \texttt{hospital} & Tagged
hospitals \\
\texttt{amenity} & \texttt{clinic} & \texttt{hospital} & Polyclinics,
walk-in clinics \\
\texttt{amenity} & \texttt{doctors} & \texttt{hospital} & General
practitioner offices \\
\texttt{healthcare} & \texttt{hospital} & \texttt{hospital} &
Healthcare-namespace hospitals \\
\texttt{healthcare} & \texttt{clinic} & \texttt{hospital} &
Healthcare-namespace clinics \\
\texttt{amenity} & \texttt{police} & \texttt{police} & Police
stations \\
\texttt{emergency} & \texttt{ambulance\_station} & \texttt{ambulance} &
Dedicated ambulance bases \\
\texttt{amenity} & \texttt{ambulance\_station} & \texttt{ambulance} &
Alternate tagging \\
\texttt{amenity} & \texttt{fire\_station} & \texttt{ambulance} & Fire
stations (often run ambulances) \\
\texttt{service:vehicle:recovery} & \texttt{yes} & \texttt{towing} &
Vehicle recovery / breakdown \\
\texttt{service:vehicle:tow} & \texttt{yes} & \texttt{towing} & Tow
truck services \\
\texttt{amenity} & \texttt{vehicle\_recovery} & \texttt{towing} &
Dedicated recovery yards \\
\texttt{shop} & \texttt{car\_repair} & \texttt{repair} & Mechanic /
garage \\
\texttt{amenity} & \texttt{car\_repair} & \texttt{repair} & Alternate
tagging \\
\texttt{shop} & \texttt{tyres} & \texttt{tyre} & Tyre / puncture
shops \\
\texttt{shop} & \texttt{tyre} & \texttt{tyre} & Alternate spelling \\
\texttt{shop} & \texttt{car} & \texttt{showroom} & Car dealerships /
showrooms \\
\texttt{shop} & \texttt{car\_parts} & \texttt{showroom} & Auto parts
dealers \\
\end{longtable}

\hypertarget{category-alignment-with-rulebook}{%
\subsubsection{2.3 Category alignment with
rulebook}\label{category-alignment-with-rulebook}}

The rulebook (Section 1.3.3 ``Key Aspects for Coders to Include'')
explicitly lists: - Nearest \textbf{Police Station} → \texttt{police}
category - \textbf{Hospitals} → \texttt{hospital} category -
\textbf{Ambulance services} → \texttt{ambulance} category -
\textbf{Towing services} → \texttt{towing} category ✅ - \textbf{Nearest
puncture shops} → \texttt{tyre} category ✅ - \textbf{Showrooms} →
\texttt{showroom} category ✅

All six rulebook-mandated categories are covered.

\begin{center}\rule{0.5\linewidth}{0.5pt}\end{center}

\hypertarget{search-result-contact-object}{%
\subsection{3. Search Result Contact
Object}\label{search-result-contact-object}}

\textbf{Source:} \texttt{backend/services/overpass\_service.py} →
\texttt{parse\_element()} \textbf{Returned by:} \texttt{GET\ /search},
\texttt{POST\ /triage}, cached in \texttt{localStorage}

\hypertarget{schema-2}{%
\subsubsection{3.1 Schema}\label{schema-2}}

\begin{verbatim}
{
  "id":         "string  — Stable unique ID, format 'osm_{id}' or 'gp_{place_id}'",
  "name":       "string  — Establishment name (English preferred, falls back to local)",
  "category":   "enum    — hospital | police | ambulance | towing | repair | tyre | showroom",
  "phone":      "string|null — E.164-normalised phone, or basic-cleaned fallback",
  "distance":   "float   — Great-circle distance from user in kilometres, 2-decimal",
  "lat":        "float   — Latitude of establishment, WGS84",
  "lon":        "float   — Longitude of establishment, WGS84",
  "source":     "string  — 'OpenStreetMap' or 'Google Places'",
  "isOpen":     "boolean|null — Parsed from opening_hours; null if unparseable",
  "aiReason":   "string|null  — Set after /triage if AI prioritised this contact"
}
\end{verbatim}

\begin{center}\rule{0.5\linewidth}{0.5pt}\end{center}

\hypertarget{triage-requestresponse}{%
\subsection{4. Triage Request/Response}\label{triage-requestresponse}}

\textbf{Endpoint:} \texttt{POST\ /triage} \textbf{Source:}
\texttt{backend/services/ai\_triage.py}

\hypertarget{request-schema}{%
\subsubsection{4.1 Request schema}\label{request-schema}}

\begin{verbatim}
{
  "injured":  "boolean — Are there injured persons on scene?",
  "blocking": "boolean — Is the vehicle blocking traffic?",
  "contacts": "Contact[] — Array of contact objects from /search"
}
\end{verbatim}

\hypertarget{response-schema}{%
\subsubsection{4.2 Response schema}\label{response-schema}}

\begin{verbatim}
{
  "contacts": "Contact[] — Same contacts, AI-reordered, top one has aiReason set",
  "reason":   "string — One-line explanation of the prioritisation"
}
\end{verbatim}

\hypertarget{priority-matrix-rule-based-fallback}{%
\subsubsection{4.3 Priority matrix (rule-based
fallback)}\label{priority-matrix-rule-based-fallback}}

\begin{longtable}[]{@{}
  >{\raggedright\arraybackslash}p{(\columnwidth - 4\tabcolsep) * \real{0.3333}}
  >{\raggedright\arraybackslash}p{(\columnwidth - 4\tabcolsep) * \real{0.3333}}
  >{\raggedright\arraybackslash}p{(\columnwidth - 4\tabcolsep) * \real{0.3333}}@{}}
\toprule\noalign{}
\begin{minipage}[b]{\linewidth}\raggedright
\texttt{injured}
\end{minipage} & \begin{minipage}[b]{\linewidth}\raggedright
\texttt{blocking}
\end{minipage} & \begin{minipage}[b]{\linewidth}\raggedright
Top priority order
\end{minipage} \\
\midrule\noalign{}
\endhead
\bottomrule\noalign{}
\endlastfoot
true & true & ambulance → hospital → police → towing → repair → tyre →
showroom \\
true & false & ambulance → hospital → police → towing → repair → tyre →
showroom \\
false & true & police → towing → repair → ambulance → hospital → tyre →
showroom \\
false & false & repair → tyre → towing → showroom → police → ambulance →
hospital \\
\end{longtable}

Within each tier, contacts are sorted by ascending distance.

\begin{center}\rule{0.5\linewidth}{0.5pt}\end{center}

\hypertarget{dispatch-summary-request}{%
\subsection{5. Dispatch Summary
Request}\label{dispatch-summary-request}}

\textbf{Endpoint:} \texttt{POST\ /dispatch-summary} \textbf{Source:}
\texttt{backend/services/dispatch\_service.py}

\hypertarget{request-schema-1}{%
\subsubsection{5.1 Request schema}\label{request-schema-1}}

\begin{verbatim}
{
  "lat":      "float — Latitude, WGS84",
  "lon":      "float — Longitude, WGS84",
  "landmark": "string|null — Reverse-geocoded landmark text",
  "injured":  "boolean — Injured persons present",
  "blocking": "boolean — Vehicle blocking traffic"
}
\end{verbatim}

\hypertarget{response-schema-1}{%
\subsubsection{5.2 Response schema}\label{response-schema-1}}

\begin{verbatim}
{
  "summary": "string — Human-readable dispatch text, ready to read aloud or paste",
  "source":  "enum   — 'ai' (Gemini-generated) or 'template' (rule-based fallback)"
}
\end{verbatim}

\begin{center}\rule{0.5\linewidth}{0.5pt}\end{center}

\hypertarget{health-check-response}{%
\subsection{6. Health Check Response}\label{health-check-response}}

\textbf{Endpoint:} \texttt{GET\ /health} \textbf{Source:}
\texttt{backend/services/health\_service.py}

\hypertarget{schema-3}{%
\subsubsection{6.1 Schema}\label{schema-3}}

\begin{verbatim}
{
  "status":         "string — 'ok' when service is up",
  "uptime_seconds": "integer — Seconds since process start",
  "configured": {
    "gemini":        "boolean — Gemini API key present",
    "google_places": "boolean — Google Places API key present"
  },
  "cache": {
    "overpass": { "size": "int", "hits": "int", "misses": "int", "hit_rate": "float" },
    "google":   { "size": "int", "hits": "int", "misses": "int", "hit_rate": "float" },
    "geocode":  { "size": "int", "hits": "int", "misses": "int", "hit_rate": "float" }
  }
}
\end{verbatim}

\begin{center}\rule{0.5\linewidth}{0.5pt}\end{center}

\hypertarget{in-memory-ttl-cache}{%
\subsection{7. In-Memory TTL Cache}\label{in-memory-ttl-cache}}

\textbf{Source:} \texttt{backend/services/cache.py} → \texttt{TTLCache}
\textbf{Not persisted to disk} --- rebuilt on each process start.

\hypertarget{instances}{%
\subsubsection{7.1 Instances}\label{instances}}

\begin{longtable}[]{@{}
  >{\raggedright\arraybackslash}p{(\columnwidth - 6\tabcolsep) * \real{0.2500}}
  >{\raggedright\arraybackslash}p{(\columnwidth - 6\tabcolsep) * \real{0.2500}}
  >{\raggedright\arraybackslash}p{(\columnwidth - 6\tabcolsep) * \real{0.2500}}
  >{\raggedright\arraybackslash}p{(\columnwidth - 6\tabcolsep) * \real{0.2500}}@{}}
\toprule\noalign{}
\begin{minipage}[b]{\linewidth}\raggedright
Cache
\end{minipage} & \begin{minipage}[b]{\linewidth}\raggedright
TTL
\end{minipage} & \begin{minipage}[b]{\linewidth}\raggedright
Max Entries
\end{minipage} & \begin{minipage}[b]{\linewidth}\raggedright
Cached payload
\end{minipage} \\
\midrule\noalign{}
\endhead
\bottomrule\noalign{}
\endlastfoot
\texttt{overpass\_cache} & 3600 s (1 hour) & 200 & List of contact
objects per (lat, lon, radius) \\
\texttt{google\_cache} & 3600 s (1 hour) & 200 & List of contact objects
per (lat, lon, radius) \\
\texttt{geocode\_cache} & 86400 s (24 hours) & 500 &
\texttt{\{landmark,\ country\_code\}} per (lat, lon) \\
\end{longtable}

\hypertarget{cache-key-format}{%
\subsubsection{7.2 Cache key format}\label{cache-key-format}}

\begin{verbatim}
{lat:.4f}_{lon:.4f}_{suffix}
\end{verbatim}

Coordinates are rounded to 4 decimal places (\textasciitilde11 m grid)
so sub-meter GPS jitter does not cause cache misses.

\begin{center}\rule{0.5\linewidth}{0.5pt}\end{center}

\hypertarget{frontend-local-cache-offlinedb.js}{%
\subsection{8. Frontend Local Cache
(offlineDB.js)}\label{frontend-local-cache-offlinedb.js}}

\textbf{Source:} \texttt{frontend/src/utils/offlineDB.js}
\textbf{Storage:} \texttt{localStorage} \textbf{Purpose:}
Last-known-good results when the user is fully offline.

\hypertarget{schema-localstorage-value-json-encoded}{%
\subsubsection{8.1 Schema (localStorage value,
JSON-encoded)}\label{schema-localstorage-value-json-encoded}}

\begin{verbatim}
{
  "key":        "string — '{lat:.2f}_{lon:.2f}' (1km grid)",
  "contacts":   "Contact[]",
  "landmark":   "string|null",
  "countryCode":"string|null",
  "timestamp":  "integer — Unix milliseconds when cached"
}
\end{verbatim}

\hypertarget{ttl}{%
\subsubsection{8.2 TTL}\label{ttl}}

24 hours. Older entries are returned with a ``stale'' indicator but not
silently discarded --- a stale cached contact is better than nothing in
an emergency.

\begin{center}\rule{0.5\linewidth}{0.5pt}\end{center}

\hypertarget{service-worker-cache-strategy}{%
\subsection{9. Service Worker Cache
Strategy}\label{service-worker-cache-strategy}}

\textbf{Source:} \texttt{frontend/public/sw.js} (Workbox)
\textbf{Storage:} Cache Storage API (managed by browser)

\begin{longtable}[]{@{}
  >{\raggedright\arraybackslash}p{(\columnwidth - 4\tabcolsep) * \real{0.3333}}
  >{\raggedright\arraybackslash}p{(\columnwidth - 4\tabcolsep) * \real{0.3333}}
  >{\raggedright\arraybackslash}p{(\columnwidth - 4\tabcolsep) * \real{0.3333}}@{}}
\toprule\noalign{}
\begin{minipage}[b]{\linewidth}\raggedright
Pattern
\end{minipage} & \begin{minipage}[b]{\linewidth}\raggedright
Strategy
\end{minipage} & \begin{minipage}[b]{\linewidth}\raggedright
Reason
\end{minipage} \\
\midrule\noalign{}
\endhead
\bottomrule\noalign{}
\endlastfoot
\texttt{/search*} & NetworkFirst (5s timeout) & Fresh data preferred,
cache is fallback \\
\texttt{/triage*} & NetworkFirst (5s timeout) & Same \\
\texttt{/offline-pack} & CacheFirst & Static seed; only updates with new
deploys \\
\texttt{/static/*} & CacheFirst & App shell --- versioned by Vite \\
\end{longtable}

\begin{center}\rule{0.5\linewidth}{0.5pt}\end{center}

\hypertarget{summary-all-persistent-structures}{%
\subsection{10. Summary --- All persistent
structures}\label{summary-all-persistent-structures}}

\begin{longtable}[]{@{}
  >{\raggedright\arraybackslash}p{(\columnwidth - 8\tabcolsep) * \real{0.2000}}
  >{\raggedright\arraybackslash}p{(\columnwidth - 8\tabcolsep) * \real{0.2000}}
  >{\raggedright\arraybackslash}p{(\columnwidth - 8\tabcolsep) * \real{0.2000}}
  >{\raggedright\arraybackslash}p{(\columnwidth - 8\tabcolsep) * \real{0.2000}}
  >{\raggedright\arraybackslash}p{(\columnwidth - 8\tabcolsep) * \real{0.2000}}@{}}
\toprule\noalign{}
\begin{minipage}[b]{\linewidth}\raggedright
\#
\end{minipage} & \begin{minipage}[b]{\linewidth}\raggedright
Artifact
\end{minipage} & \begin{minipage}[b]{\linewidth}\raggedright
Type
\end{minipage} & \begin{minipage}[b]{\linewidth}\raggedright
Size
\end{minipage} & \begin{minipage}[b]{\linewidth}\raggedright
Update cadence
\end{minipage} \\
\midrule\noalign{}
\endhead
\bottomrule\noalign{}
\endlastfoot
1 & Emergency numbers DB & Static JSON & 200 entries &
\textasciitilde once per decade \\
2 & OSM category map & Source code & 18 mappings & When OSM tags
evolve \\
3 & Contact object & Runtime schema & Per query & Live \\
4 & Triage I/O & Runtime schema & Per query & Live \\
5 & Dispatch I/O & Runtime schema & Per query & Live \\
6 & Health response & Runtime schema & Per ping & Live \\
7 & Server TTL cache & In-memory & Up to 900 entries & Auto-evicting \\
8 & Browser localStorage & Per device & 1 entry per area & 7-day TTL \\
9 & Service Worker cache & Per device & Up to 50 entries &
NetworkFirst \\
\end{longtable}

\textbf{RoadSOS does not store any user data on a server.} No accounts,
no telemetry, no user identifiers. All user state is local to the
device.
